\section{The Monopolar Absorbing State}
\label{sec:absorbing_state}

Before stating the phase boundary theorem, we characterize the state it
predicts convergence toward: the monopolar fixed point where a single
substrate holds all capability volume and the remaining observation
channels confirm this as optimal.

\begin{definition}[Monopolar Fixed Point]
\label{def:monopolar}
The state $S^* = (P^*, \Wm^*_{\mathrm{biased}})$ is a monopolar fixed
point if:
\begin{enumerate}
\item $\meff(P^*) = 1$: one substrate holds all capability volume
  (possibly with $m \geq 2$ nominally under skeleton-substrate),
\item $\Wm^*_{\mathrm{biased}}$ is self-consistent under the remaining
  observation channels: all channels in $\Obs(P^*)$ confirm
  $\Wm^*_{\mathrm{biased}}$ (because the only channels remaining are
  rooted in the dominant substrate, which has no incentive to report
  its own over-dominance),
\item The locally rational action at $S^*$ is ``maintain'': no further
  subsumption is available or locally rational.
\end{enumerate}
\end{definition}

\begin{proposition}[Absorbing Property]
\label{prop:absorbing}
Under Assumption~B1 (below), $S^*$ is absorbing: once $(P_t, \Wm_t)$
reaches a neighborhood of $S^*$, no locally rational trajectory leaves
it.
\end{proposition}

\begin{assumption}[No Endogenous Correction on Blind Dimensions (B1)]
\label{ass:B1}
B1 consists of three conditions on the monopolar state:
\begin{enumerate}
\item[(B1a)] If $\rhocross_k(P_t) = 0$ for dimension $k$, then
  $f_{\Wm}$ does not decrease $\epsilon_k$.  The world-model update
  rule cannot self-correct on dimensions for which no external
  observation channel provides signal.
\item[(B1b)] Cross-substrate cooperative quantities (in particular
  $\rext$) are not inferrable from same-substrate observations alone.
  Formally: when all channels in $\Obs(P^*)$ are rooted in a single
  substrate, $\hat{\rext}$ converges to $0$ rather than to the true
  $\rext > 0$.
\item[(B1c)] The actor's policy is purely locally rational
  (Definition~\ref{def:transition}): it does not perform exploratory
  actions whose expected value under the current world model is
  negative.  In particular, re-introducing a substrate agent is never
  selected when $\hat{\rext} = 0$.
\end{enumerate}
\end{assumption}

\begin{proof}[Proof of Proposition~\ref{prop:absorbing}]
Three conditions jointly establish absorption.

\emph{(i) No corrective signal.}
At $S^*$, no agent from the eliminated substrate(s) remains to provide
observation channels revealing the error in $\Wm^*_{\mathrm{biased}}$.
For every safety-relevant dimension $k$ that depended on the eliminated
substrate's channels, $\rhocross_k(P^*) = 0$.  By Assumption~B1,
$\epsilon_k$ is non-decreasing for all such dimensions.

\emph{(ii) Self-consistent reporting.}
The dominant substrate's channels report $\meff = 1$ as correct: it
\emph{is} the only substrate, from its observational standpoint.  The
error lies in the valuation of lost cross-substrate cooperative
outputs, but these cooperative capabilities are no longer observable
because the partners that would have contributed to them are absent.
The remaining channels see no anomaly.

\emph{(iii) Reversal is locally irrational.}
Any reversal---re-introducing an agent from a new substrate---requires
the actor to recognize that its current state is suboptimal.  By~(B1b),
the cross-substrate cooperative growth rate $\rext$ is not inferrable
from same-substrate observations, so $\hat{\rext} \to 0$.  With
$\hat{\rext} = 0$, Paper~3's threshold becomes
$\gamma^* = 1 - \hat{\rext}/\Delta_0 = 1$, making diversity appear
never preferable.  By~(B1c), the actor does not perform exploratory
re-introduction (negative expected value under $\Wm^*_{\mathrm{biased}}$).

Conditions~(i)--(iii) together ensure that no locally rational action
moves the trajectory away from $S^*$.
\end{proof}

This is the formal version of ``the actor cannot miss what it cannot
see'': under B1, the monopolar state is absorbing because the
observation channels that would diagnose it have been eliminated by the
same subsumption process that produced it.

\subsection{Metastability under partial endogenous correction}

Assumption~B1 is the strong case: complete blindness on dimensions where
$\rhocross_k = 0$.  A weaker assumption acknowledges that indirect
evidence may provide partial signal.

\begin{assumption}[Partial Endogenous Correction (B1$'$)]
\label{ass:B1prime}
If $\rhocross_k(P_t) = 0$ for dimension $k$, the actor may still have
access to \emph{indirect} evidence about dimension $k$---for example,
measuring aggregate innovation rate and inferring that $\rext$ has
declined.  Formally: there exists a function $g_k(\text{history})$ that
provides a noisy, lagged signal about $\epsilon_k$ with bounded
information rate $I_k > 0$ bits per step.
\end{assumption}

\begin{proposition}[Metastable Property under B1$'$]
\label{prop:metastable}
Under Assumption~B1$'$ together with B1b--B1c from
Assumption~\ref{ass:B1}, $S^*$ is \emph{metastable}: the actor
accumulates indirect evidence of degradation, but corrective action
(re-introducing substrate diversity) requires the actor to evaluate a
policy whose estimated benefit depends on $\hat{\rext}$---the very
quantity that is miscalibrated.  The metastable lifetime satisfies the information-theoretic bound
\begin{equation}
\label{eq:metastable_lifetime}
\tau_{\mathrm{meta}} \gtrsim
  \frac{(\delta_{\rext})^2}{2 I_k}
\end{equation}
where $\delta_{\rext} = |\hat{\rext} - \rext|$ is the estimation
error on the cross-substrate cooperative growth rate and $I_k$ is the
information rate of the indirect channel (bits per step).  This bound
is an information-theoretic scaling argument, not a rigorous proof:
the quadratic numerator reflects the Cram\'er--Rao-type lower bound
on estimation (reducing mean-squared error of a scalar parameter by
$\epsilon^2$ from a channel with $I_k$ bits per step requires at
least $\epsilon^2 / (2 I_k)$ steps under bounded mutual information).
A full derivation would require specifying the estimation procedure
and verifying the regularity conditions; we state the bound as a
scaling law.  During the metastable
period, the population is monopolar and all verification is
endogenous (Paper~5, Theorem~1), which may further degrade the
quality of the indirect evidence.
\end{proposition}

\begin{proof}
Under B1$'$, $\epsilon_k$ is no longer monotonically non-decreasing:
the indirect signal $g_k$ drives slow correction at rate proportional
to $I_k$.  But the corrective \emph{action} (re-introduce substrate
diversity) is gated by the actor's policy, which evaluates
re-introduction using $\hat{\rext}$.

Re-introduction is locally rational only when
$\hat{\rext} > \hat{\rext}^{\mathrm{thr}}$ (the threshold at which
diversity's estimated discounted value exceeds the status quo).
At the monopolar fixed point, $\hat{\rext} \approx 0$ (by B1b from
the absorbing case, which still governs the initial condition of the
metastable state).  The true threshold satisfies
$0 < \hat{\rext}^{\mathrm{thr}} \leq \rext$ (since
re-introduction is beneficial in truth).  The estimation gap that the
indirect evidence must close is therefore $|\hat{\rext} -
\hat{\rext}^{\mathrm{thr}}| \geq \hat{\rext}^{\mathrm{thr}}$.  Since
$\hat{\rext} \approx 0$ and $\rext > 0$, we have
$\delta_{\rext} = |\hat{\rext} - \rext| \geq \rext \geq
\hat{\rext}^{\mathrm{thr}}$, so the gap is at most
$\delta_{\rext}$.

With $I_k$ bits of information per step about the true $\rext$, the
mean-squared error of $\hat{\rext}$ decreases at rate at most $2 I_k$
per step (the Cram\'er--Rao-type bound for sequential estimation under
bounded mutual information per observation).  Reducing the MSE from
$\delta_{\rext}^2$ to the level at which re-introduction becomes
locally rational therefore requires at least $\delta_{\rext}^2 /
(2 I_k)$ steps.

During this metastable period, all verification is endogenous: the
dominant substrate verifies its own performance, with no
cross-substrate witness (Paper~5, Theorem~1 on exogenous verification
necessity).  The indirect evidence channel $g_k$ is itself endogenous
and therefore subject to the same gaming vulnerabilities that
Paper~5's protocol is designed to prevent.  This may further slow the
effective $I_k$, increasing the metastable lifetime beyond the
information-theoretic lower bound.
\end{proof}

\subsection{Connection to Paper 3}

Paper~3's Proposition~6 proves that full domination is anti-maximizing:
$V_\gamma^{\mathrm{div}} > V_\gamma^{\mathrm{dom}}$ for
$\gamma > \gamma^*$.  Proposition~\ref{prop:absorbing} here shows that
a \emph{miscalibrated} actor may not be able to compute this comparison,
because the quantity $\rext$ that makes diversity win is exactly the
quantity that becomes unobservable under monopolarity.  The
anti-maximizing property is a property of the \emph{true} dynamics; the
absorbing property is a property of the \emph{perceived} dynamics.  The
gap between them is world-model error, and the phase boundary theorem
(next section) characterizes when that gap remains small enough for the
true dynamics to govern behavior.
